\documentclass[oneside,a4paper,14pt]{extarticle}
\usepackage[a4paper,letterpaper,top=20mm,bottom=20mm,left=20mm,right=10mm]{geometry}
\usepackage[russian]{babel}
\usepackage{indentfirst}
\usepackage{graphicx}
\usepackage{caption}
\usepackage{titlesec}
\usepackage{minted, fancyvrb}
\usepackage{hyperref}
\usepackage{enumitem}
\usepackage{amsmath}
\usepackage{fancyhdr}
\usepackage{colortbl}
\usepackage[dvipsnames]{xcolor}
\usepackage{tikz-timing}
\usepackage{tikz}
\usetikzlibrary{karnaugh}
\usepackage{float}

% Настройка автоматических подписей к рисункам (Замена Рис. на Рисунок и тире в качестве разделителя)
\DeclareCaptionLabelSeparator{customdash}{~---~}
\captionsetup[figure]{name=Рисунок, labelsep=customdash, justification=centering}

\setminted{style = rainbow_dash, fontsize = \small}

% Форматирование заголовков
\titleformat{\section}{\normalsize\bfseries}{\thesection}{1em}{}
\titleformat{\subsection}{\normalsize\bfseries}{\thesubsection}{1em}{}
\titleformat{\subsubsection}{\normalsize\bfseries}{\thesubsubsection}{1em}{}

% Интерлиньяж и абзац
\renewcommand\baselinestretch{1.33}
\setlength{\parindent}{1.25cm} % длина красной строки

% Для всех списков
\setlist[enumerate]{
  left=0.5\parindent,
  label=\arabic*.,
  itemsep=0pt,
  topsep=5pt,
  partopsep=0pt,
  parsep=0pt
}
\setlist[itemize]{
  left=0.5\parindent,
  itemsep=0pt,
  topsep=5pt,
  partopsep=0pt,
  parsep=0pt
}

% Гиперссылки
\hypersetup{
  colorlinks=true,
  linkcolor=black,
  urlcolor=blue,
  pdfborder={0 0 0},
  pdftitle={Контрольная работа: Построение 2-разрядного сдвигового регистра на JK-триггерах},
  pdfauthor={Черкасов А.А. и Зинин В.А.},
}

\begin{document}

\newpage
\thispagestyle{empty}
\begin{center}
  МИНИСТЕРСТВО НАУКИ И ВЫСШЕГО ОБРАЗОВАНИЯ РОССИЙСКОЙ ФЕДЕРАЦИИ\\
  ФЕДЕРАЛЬНОЕ ГОСУДАРСТВЕННОЕ БЮДЖЕТНОЕ УЧРЕЖДЕНИЕ ВЫСШЕГО ОБРАЗОВАНИЯ\\
  <<ВЯТСКИЙ ГОСУДАРСТВЕННЫЙ УНИВЕРСИТЕТ>>\\
  Институт математики и информационных систем\\
  Факультет автоматики и вычислительной техники\\
  Кафедра электронных вычислительных машин
\end{center}
\vspace{10mm}

\hfill
\begin{tabular}{l}
  \footnotesize Дата сдачи на проверку:                                          \\
  \footnotesize <<\rule[-1mm]{5mm}{0.10mm}\/>>\rule[-1mm]{20mm}{0.10mm}\ 2026 г. \\
  \footnotesize Проверено:                                                       \\
  \footnotesize <<\rule[-1mm]{5mm}{0.10mm}\/>>\rule[-1mm]{20mm}{0.10mm}\ 2026 г. \\
\end{tabular}
\vfill

\begin{center}
  Отчёт по контрольной работе\\
  по теме\\
  <<Построение 2-разрядного сдвигового регистра на JK-триггерах>>\\
  по дисциплине\\
  <<Цифровая схемотехника>>\\
\end{center}
\vspace{25mm}
\noindent
\begin{tabular}{ll}
  Разработали студенты гр. ИВТб-2301-05-00 & \hspace{5mm}\rule[-1mm]{30mm}{0.10mm}\,/Черкасов А. А./ \\
                                           & \hspace{5mm}\rule[-1mm]{30mm}{0.10mm}\,/Зинин В. А./    \\
                                           & \hspace{12.5mm}\footnotesize(подпись)                   \\
  Преподаватель                            & \hspace{5mm}\rule[-1mm]{30mm}{0.10mm}\,/Мельцов В. Ю./  \\
                                           & \hspace{12.5mm}\footnotesize(подпись)                   \\
\end{tabular}

\noindent
\begin{tabular}{lp{58mm}r}
  Работа защищена &  & \hspace{5mm}<<\rule[-1mm]{5mm}{0.10mm}\/>>\rule[-1mm]{30mm}{0.10mm}\ 2026 г.
\end{tabular}
\vfill

\begin{center}
  Киров\\
  2026
\end{center}

\newpage\thispagestyle{plain}

\section{Цель работы}
Получить теоретические и практические навыки проектирования и анализа последовательностных цифровых устройств. Синтезировать функциональную и принципиальную схемы 2-разрядного универсального сдвигового регистра на базе универсальных JK-триггеров с использованием интегральных микросхем серий КП7 и ТВ6 (ТТЛ-логика).

\section{Теоретическое описание}
\subsection{JK-триггер}
JK-триггер является универсальным двухступенчатым автоматом памяти, имеющим информационные входы $J$ (Jump --- установка <<1>>) и $K$ (Kill --- сброс в <<0>>), тактовый синхронизирующий вход $C$, а также независимые асинхронные входы инверсной установки $\overline{S}$ и сброса $\overline{R}$. Особенностью JK-триггера является исключение неопределенного состояния: при подаче комбинации $J = K = 1$ автомат переходит в режим инверсии предыдущего состояния (счётный режим).

\begin{center}
  Таблица 1 --- Таблица истинности JK-триггера \\
  \vspace{3mm}
  \begin{tabular}{|c|c|c|c|c|c|}
    \hline
    $\overline{S}$ & $\overline{R}$ & C & J & K & $Q_{n+1}$ \\
    \hline
    0 & 1 & X & X & X & 1 \\
    1 & 0 & X & X & X & 0 \\
    0 & 0 & X & X & X & 1* \\
    1 & 1 & $\downarrow$ & 0 & 0 & $Q_n$ \\
    1 & 1 & $\downarrow$ & 1 & 0 & 1 \\
    1 & 1 & $\downarrow$ & 0 & 1 & 0 \\
    1 & 1 & $\downarrow$ & 1 & 1 & $\overline{Q}_n$ \\
    \hline
  \end{tabular}
  \vspace{2mm}
  \par\footnotesize * --- неустойчивое (запрещенное) состояние при асинхронном управлении.
\end{center}

\subsection{Универсальный сдвиговый регистр}
Сдвиговый регистр представляет собой последовательностный цифровой узел, осуществляющий направленный сдвиг хранимого многоразрядного слова под воздействием тактовых импульсов. В проектируемом 2-разрядном универсальном сдвиговом регистре реализована многофункциональная структура, управляемая кодовыми комбинациями на селекторных входах $S_1$ и $S_0$:
\begin{itemize}
  \item $S_1S_0 = 00$ --- режим хранения данных (состояние выходов стабильно);
  \item $S_1S_0 = 01$ --- последовательный сдвиг влево (информационный вход $DL$);
  \item $S_1S_0 = 10$ --- последовательный сдвиг вправо (информационный вход $DR$);
  \item $S_1S_0 = 11$ --- параллельная асинхронная или синхронная загрузка слова ($D_1$ --- старший бит, $D_0$ --- младший бит).
\end{itemize}
Синхронизация всех разрядов осуществляется единой шиной $CLK$, объединяющей динамические тактовые входы триггеров.

\begin{center}
  Таблица 2 --- Режимы работы универсального сдвигового регистра \\
  \vspace{3mm}
  \begin{tabular}{|c|c|c|}
    \hline
    $S_1$ & $S_0$ & Режим функционирования \\
    \hline
    0 & 0 & Хранение информации \\
    0 & 1 & Сдвиг в левую сторону ($DL$) \\
    1 & 0 & Сдвиг в правую сторону ($DR$) \\
    1 & 1 & Параллельный ввод данных ($D_1, D_0$) \\
    \hline
  \end{tabular}
\end{center}

\section{Функциональная схема универсального сдвигового регистра}
На основании алгоритма функционирования разработана каноническая функциональная схема многофункционального регистра на JK-триггерах. Устройство содержит внешние линии последовательного ввода $DL$ и $DR$, шину параллельной загрузки данных $D_1$ и $D_0$, управляющий селектор режимов $S_0, S_1$ и тактовый вход $CLK$. Текущие логические состояния разрядов снимаются с прямых выходов триггеров $Q_0$ и $Q_1$.

\begin{figure}[H]
  \centering
  \includegraphics[width=0.9\textwidth]{pics/functional_diagram.png}
  \caption{Функциональная схема 2-разрядного универсального сдвигового регистра}
\end{figure}

\section{Принципиальная схема устройства}
Схема электрическая принципиальная спроектирована на элементной базе транзисторно-транзисторной логики (ТТЛ). В качестве элементной базы применены интегральные микросхемы \textbf{КП7} (селекторы-мультиплексоры для коммутации режимов) и \textbf{ТВ6} (универсальные двухступенчатые JK-триггеры с логическим объединением по входам). Подключение выводов корпусов ИС выполнено в соответствии с их заводскими цоколёвками, шина питания рассчитана на стандартное напряжение $+5$~В.

\begin{figure}[H]
  \centering
  \includegraphics[width=0.95\textwidth]{pics/principle_diagram.png}
  \caption{Принципиальная электрическая схема регистра на ИС КП7 и ТВ6}
\end{figure}

\section*{Вывод}
В ходе выполнения контрольной работы были изучены методы структурного синтеза последовательностных цифровых автоматов на примере 2-разрядного универсального сдвигового регистра. Спроектированное устройство успешно реализует четыре базовых микрооперации: хранение, двунаправленный сдвиг (влево/вправо) и параллельный занос данных. Была освоена методика применения ИС серий КП7 и ТВ6 ТТЛ-логики для схемотехнического воплощения спроектированной структуры с учетом обеспечения требуемых режимов синхронизации.

\end{document}